%------------------------------------------------------------------------------
\chapter{0-D Discretization}\label{sec:0d_discretization}
%------------------------------------------------------------------------------
\section{0-D Source/sink term }\label{sec:0d_source_and_sink}
In this section a zero-dimensional model is implementation of the source term, representing the source and sink of external influences, like a power plant.
The main (simple) equation will look like:
\begin{align}
    \pdiff{\vec{u}}{t} = \vec{f}(\vec{u},t)
\end{align}
Here we will show the mathematical and implementation of an air pollution model, \autoref{sec:air_pollution} \citep[ex.\ 1.1 pg.\ 7]{HundsdorferAndVerwer2003} and a Brusselator, \autoref{sec:brusselator}  \citep{AultHolmgreen2003}.
Some numerical results are shown in \autoref{sec:0d_numerical_experiments}.
%------------------------------------------------------------------------------
\subsection{Air pollution}\label{sec:air_pollution}
%------------------------------------------------------------------------------
\subsubsection{Analytic description}
We illustrate the mass action law by the following three reactions between oxygen $O_2$, atomic oxygen $O$, nitrogen oxide $NO$, and nitrogen dioxide $NO_2$ \citep[eq.\ 1.1, page 7]{HundsdorferAndVerwer2003}:
\begin{align}
    NO_2 + h\nu & \xrightarrow{k_1} NO + O, \\
    O + O_2 & \xrightarrow{k_2} O_3, \\
    NO + O_3 & \xrightarrow{k_3} O_2 + NO_2.
\end{align}
These reactions are basic to any tropospheric air pollution model.
The first reaction is photochemical and says that $NO$ and $O$ are formed from $NO_2$ by
photo-dissociation caused by solar radiation, indicated by $h\nu$.
This depends on the time of the day and therefore $k_1 = k_1(t)$.
We consider the concentrations $u_1 = \bunit{O}$, $u_2 = \bunit{NO}$, $u_3 = \bunit{NO2}$, $u_4 = \bunit{O3}$.
The oxygen concentration $\bunit{O_2}$ is treated as constant, and we assume there is a constant source term $\sigma_2$ simulating emission of nitrogen oxide.
The rate functions and stoichiometric matrix are
\begin{align}
    g(t,u) =
    \begin{pmatrix}
        k_1(t) u_3 \\
        k_2u_1 \\
        k_3u_2u_4
    \end{pmatrix}
    ,
    \qquad
    S =
    \begin{pmatrix}
        1 & -1 & 0 \\
        1 & 0 & -1 \\
        -1 & 0 & 1 \\
        0 & 1 & -1
    \end{pmatrix}
\end{align}
The corresponding ODE system reads:
\begin{align}
    \pdiff{u_1}{t} & = k_1 u_3 -k_2 u_1
    \\
    \pdiff{u_2}{t} & = k_1 u_3 - k_3 u_2 u_4 +\sigma_2
    \\
    \pdiff{u_3}{t} & = k_3 u_2 u_4 - k_1 u_3
    \\
    \pdiff{u_4}{t} & = k_2 u_1 - k_3 u_2 u_4
\end{align}
We see that
\begin{align}
    \pdiff{u_1}{t} + \pdiff{u_3}{t} + \pdiff{u_4}{t} = 0,
\end{align}
and
\begin{align}
    \pdiff{u_2}{t} + \pdiff{u_3}{t} = \sigma_2 t
\end{align}
hence $\bunit{O} + \bunit{NO2} + \bunit{O3}$ is a conserved quantity while $\bunit{NO} + \bunit{NO2}$ growths
with $\sigma_2t$.
By considering the null-space of $S^T$ it is seen that these are the two
mass laws for this chemical model.

With $\vec{u}(0) = (0.0, \num{2.0e-1}, \num{2.0e-3}, \num{2.0e-1})^T$ and $\sigma_2 = \num{e-7}$, and the coefficients $k$ are defined as:
\begin{align}
    k_1 & = \begin{cases}
        10^{-5} \exp{(7\ {\mathit g}(t))}
        \\
        \num{e-40}, \qquad \text{during  night}
    \end{cases}
    \\
    k_2 & = \num{2.0e-2}
    \\
    k_3 & = \num{1.0e-3}
\end{align}
with
\begin{align}
    {\mathit g}(t) =\left(\sin\left(\frac{\pi}{16} (t_h - 4)\right)\right)^{0.2}, \qquad t_h = \frac{t}{3600};
\end{align}
where $t_h$ is the time in hours.
The given conditions are different from the conditions as defined in \citet[pg.\ 8]{HundsdorferAndVerwer2003}
How these equations are discretized is given in \autoref{sec:air_pollution_discretization}.
%
%-------------------------------------------------------------------------------
\subsubsection{Numerical discretization}\label{sec:air_pollution_discretization}
The discretization in $\Delta$-formulation reads:
\begin{align}
    \frac{1}{\Dt}\Delta u_1^{n+1, p+1} & = - \frac{1}{\Dt}(u_1^{n+1,p} - u_1^n) + k_1 (u_3^{n+\theta,p+1}) - k_2 (u_1^{n+\theta,p+1})
    \\
    \frac{1}{\Dt}\Delta u_2^{n+1, p+1} & = - \frac{1}{\Dt}(u_2^{n+1,p} - u_2^n) + k_1(u_3^{n+\theta,p+1}) - k_3 (u_2^{n+\theta,p+1}) (u_4^{n+\theta,p+1}) +\sigma_2
    \\
    \frac{1}{\Dt}\Delta u_3^{n+1, p+1} & = - \frac{1}{\Dt}(u_3^{n+1,p} - u_3^n) + k_3 (u_2^{n+\theta,p+1}) (u_4^{n+\theta,p+1}) - k_1(u_3^{n+\theta,p+1})
    \\
    \frac{1}{\Dt}\Delta u_4^{n+1, p+1} & = - \frac{1}{\Dt}(u_4^{n+1,p} - u_4^n) + k_2(u_1^{n+\theta,p+1}) - k_3 (u_2^{n+\theta,p+1}) (u_4^{n+\theta,p+1})
\end{align}
with linearization of $\vec{u}^{n+\theta,p+1}$ yields:
\begin{align}
    \frac{1}{\Dt}\Delta u_1^{n+1, p+1} & =
    - \frac{1}{\Dt}(u_1^{n+1,p} - u_1^n) +
    \nonumber \\*
    & + k_1 (u_3^{n+\theta,p} + \theta \Delta u_3^{n+1, p+1}) - k_2 (u_1^{n+\theta,p} + \theta \Delta u_1^{n+1, p+1})
    \\
    %-----
    \frac{1}{\Dt}\Delta u_2^{n+1, p+1} & = - \frac{1}{\Dt}(u_2^{n+1,p} - u_2^n) + k_1(u_3^{n+\theta,p} + \theta \Delta u_3^{n+1, p+1}) +
    \nonumber \\*
    & - k_3 (u_2^{n+\theta,p} + \theta \Delta u_2^{n+1, p+1}) (u_4^{n+\theta,p} + \theta \Delta u_4^{n+1, p+1}) +\sigma_2
    \\
    %-----
    \frac{1}{\Dt}\Delta u_3^{n+1, p+1} & = - \frac{1}{\Dt}(u_3^{n+1,p} - u_3^n) + k_3 (u_2^{n+\theta,p} + \theta \Delta u_2^{n+1, p+1}) (u_4^{n+\theta,p} +
    \nonumber \\*
    & +  \theta \Delta u_4^{n+1, p+1}) - k_1(u_3^{n+\theta,p} + \theta \Delta u_3^{n+1, p+1})
    \\
    %-----
    \frac{1}{\Dt}\Delta u_4^{n+1, p+1} & = - \frac{1}{\Dt}(u_4^{n+1,p} - u_4^n) + k_2(u_1^{n+\theta,p} + \theta \Delta u_1^{n+1, p+1}) +
    \nonumber \\*
    & - k_3 (u_2^{n+\theta,p} + \theta \Delta u_2^{n+1, p+1}) (u_4^{n+\theta,p} + \theta \Delta u_4^{n+1, p+1})
\end{align}
and rearrange the system of equations to $\mat{A}\vec{x}=\vec{b}$, yields
\begin{align}
    \frac{1}{\Dt}\Delta u_1^{n+1, p+1} &  - k_1  \theta \Delta u_3^{n+1, p+1} + k_2  \theta \Delta u_1^{n+1, p+1}  =
    \nonumber \\*
    & = - \frac{1}{\Dt}(u_1^{n+1,p} - u_1^n) + k_1 u_3^{n+\theta,p} - k_2 u_1^{n+\theta,p}
    \\
    %-----
    \frac{1}{\Dt}\Delta u_2^{n+1, p+1} & - k_1 \theta \Delta u_3^{n+1, p+1}
    + k_3 \theta u_4^{n+1,p} \Delta u_2^{n+1, p+1} + k_3 \theta u_2^{n+1,p} \Delta u_4^{n+1, p+1}  =
    \nonumber \\*
    & = - \frac{1}{\Dt}(u_2^{n+1,p} - u_2^n) + k_1 u_3^{n+\theta,p} - k_3 u_2^{n+\theta,p} u_4^{n+\theta,p} +\sigma_2
    \\
    %-----
    \frac{1}{\Dt}\Delta u_3^{n+1, p+1} & - k_3 u_2^{n+\theta,p} \theta \Delta u_4^{n+1, p+1} - k_3 u_4^{n+\theta,p} \theta \Delta u_2^{n+1, p+1} + k_1 \theta \Delta u_3^{n+1, p+1}  =
    \nonumber \\*
    & = - \frac{1}{\Dt}(u_3^{n+1,p} - u_3^n) + k_3 u_2^{n+\theta,p} u_4^{n+\theta,p} - k_1 u_3^{n+\theta,p}
    \\
    \frac{1}{\Dt}\Delta u_4^{n+1, p+1}&  - k_2 \theta \Delta u_1^{n+1, p+1}  + k_3 u_2^{n+\theta,p}\theta \Delta u_4^{n+1, p+1} +k_3 u_4^{n+\theta,p} \theta \Delta u_2^{n+1, p+1} =
    \nonumber \\*
    & = - \frac{1}{\Dt}(u_4^{n+1,p} - u_4^n) + k_2 u_1^{n+\theta,p} - k_3 u_2^{n+\theta,p} u_4^{n+\theta,p}
\end{align}
%------------------------------------------------------------------------------
\subsection{Brusselator}\label{sec:brusselator}
%------------------------------------------------------------------------------
\subsubsection{Analytic description}
The  ODE system for the Brusselator reads \citet[eq.\ 14,15]{AultHolmgreen2003}:
\begin{subequations}
\begin{align}
    \pdiff{u_1}{t} & = 1 - (k_2 + 1) u_1 + k_1 u_1^2 u_2,
    \\
    \pdiff{u_2}{t} & = k_2 u_1 - k_1 u_1^2 u_2
\end{align}
\label{eq:brusselator}
\end{subequations}
with $k_1 =1$ and  $k_2 = 2.5$ and initial values  $u_1(0)=0$ and $u_2(0) = 0$.
%-------------------------------------------------------------------------------
\subsubsection{Numerical discretization}\label{sec:brusselator_discretization}

The  ODE system for the Brusselator reads \citep[eq.\ 14,15]{AultHolmgreen2003}:
\begin{align}
    \pdiff{u_1}{t} & = 1 - (k_2 + 1) u_1 + k_1 u_1^2 u_2,
    \\
    \pdiff{u_2}{t} & = k_2 u_1 - k_1 u_1^2 u_2
\end{align}
The discretization in $\Delta$-formulation reads:
\begin{align}
    \frac{1}{\Dt}\Delta u_1^{n+1, p+1} & = - \frac{1}{\Dt}(u_1^{n+1,p} - u_1^n) + 1 - (k_2 +1) u_1^{n+\theta,p+1} +
    \nonumber \\*
    &  + k_1 \left(u_1^{n+\theta,p+1}\right)^2 u_2^{n+\theta,p+1}
    \\
    \frac{1}{\Dt}\Delta u_2^{n+1, p+1} & = - \frac{1}{\Dt}(u_2^{n+1,p} - u_2^n) + k_2(u_1^{n+\theta,p+1})  +
    \nonumber \\*
    & - k_1 \left(u_1^{n+\theta,p+1}\right)^2 u_2^{n+\theta,p+1}
\end{align}
with linearization of $\vec{u}^{n+\theta,p+1}$ yields:
\begin{align}
    \frac{1}{\Dt}\Delta u_1^{n+1, p+1} & =
    - \frac{1}{\Dt}(u_1^{n+1,p} - u_1^n) + 1 - (k_2 +1) \left(u_1^{n+\theta,p} + \theta \Delta u_1^{n+1, p+1}\right)+
    \nonumber \\*
    &  + k_1 \left(u_1^{n+\theta,p}+\Delta u_1^{n+1, p+1}\right)^2 \left(u_2^{n+\theta,p}+ \Delta u_1^{n+1, p+1}\right)
    \\
    \frac{1}{\Dt}\Delta u_2^{n+1, p+1} & = - \frac{1}{\Dt}(u_2^{n+1,p} - u_2^n) + k_2 (u_1^{n+\theta,p} + \theta \Delta u_1^{n+1, p+1}) +
    \nonumber \\*
    &
    - k_1 \left(u_1^{n+\theta,p}+\Delta u_1^{n+1, p+1}\right)^2 \left(u_2^{n+\theta,p}+ \Delta u_1^{n+1, p+1}\right)
\end{align}
and rearrange the system of equations to $\mat{A}\vec{x}=\vec{b}$ and omitting the second order terms, yields
\begin{align}
    &\frac{1}{\Dt}\Delta u_1^{n+1, p+1}
    - \theta \left(\left(k_2 +1\right) + 2  k_1 u_1^{n+\theta,p} u_2^{n+\theta,p}\right)\Delta u_1 ^{n+1, p+1}
    - \theta k_1 \left(u_1^{n+1, p+1}\right)^2 \Delta u_2 ^{n+1, p+1} =
    \nonumber \\*
    & = - \frac{1}{\Dt}(u_1^{n+1,p} - u_1^n) + 1 - (k_2 +1) u_1^{n+\theta,p} + k_1 \left(u_1^{n+\theta,p}\right)^2 u_2^{n+\theta,p}
    \\
    & \frac{1}{\Dt}\Delta u_2^{n+1, p+1}
    + \theta \left( k_2 + 2 k_1 u_1^{n+\theta,p} u_2^{n+\theta,p} \right) \Delta u_1 ^{n+1, p+1}
    + \theta k_1\left(u_1^{n+1, p+1}\right)^2 \Delta u_2 ^{n+1, p+1}
    = \nonumber \\*
    & =- \frac{1}{\Dt}(u_2^{n+1,p} - u_2^n) + k_2 u_1^{n+\theta,p}
    - k_1 \left(u_1^{n+\theta,p}\right)^2 u_2^{n+\theta,p}
\end{align}
This system can be implemented and solved, some results are presented in \autoref{sec:num_exp_brusselator}.
